\section{Appendix}

\subsection{Training Hardware}
\label{sec:training_hardware}

Training was performed on a desktop system with the specifications outlined in
Table~\ref{tab:training_hardware}.

\begin{table*}[!htb]
	\centering
	\small
	\begin{tabularx}{\textwidth}{@{\extracolsep{\fill}}>{\raggedright\arraybackslash}X >{\raggedright\arraybackslash}X}
		\toprule
		\textbf{Component} & \textbf{Specification}                         \\
		\midrule
		CPU                & AMD Ryzen 9 7800X3D 8-core 16-thread processor \\
		\addlinespace
		GPU                & NVIDIA GeForce RTX 4070 Super 12GB             \\
		\addlinespace
		RAM                & 32GB DDR5 6000MT/s                             \\
		\addlinespace
		Storage            & 1TB NVMe SSD                                   \\
		\addlinespace
		Operating System   & Windows 11                                     \\
		\bottomrule
	\end{tabularx}
	\caption{Training hardware specifications.}
	\label{tab:training_hardware}
\end{table*}

\subsection{Training Configuration}
\label{sec:config}

The training and export configuration for the primary detector model is shown
in listing~\ref{lst:config}.

\lstinputlisting[
	language=json,
	caption={YOLO training configuration file.},
	label={lst:config}
]{../training/source/config.json}

\subsection{Pipeline Configuration}

The configuration of the DeepStream pipeline is defined programmatically in
lines 10-19 of \texttt{builder.py} shown in listing~\ref{lst:pipeline}.

\lstinputlisting[
	style=pythonstyle,
	caption={DeepStream pipeline configuration.},
	label={lst:pipeline},
	firstline=7,
	lastline=19
]{../inference/source/pipeline/builder.py}

\subsection{Configuration for Primary Detector}

The configuration file for the primary inference model is shown in
listing~\ref{lst:primary_config}.

\lstinputlisting[
	language=deepstream,
	caption={Primary detector configuration.},
	label={lst:primary_config}
]{../inference/config/config_infer_primary.txt}

\subsection{Example Ground Truth Annotation}

An example of the ground truth annotation format is shown in
listing~\ref{lst:annotation_format}. The script for producing these annotations
is available in the GitHub repository,
\url{https://github.com/maxnaro/ABODA/blob/master/source/annotate.py}, a fork
of the original ABODA repository.

\begin{lstlisting}[
	language=json,
	caption={Ground truth annotation format example.},
	label={lst:annotation_format}
]
{
  "video1.avi": [
    {
      "has_abandonment": true,
      "true_abandon_frame": 2016,
      "bag_roi": [
        133,
        323,
        54,
        46
      ],
      "radius_px": 200,
      "threshold_s": 15
    }
  ],
  "video2.avi": [
    {
      "has_abandonment": true,
      "true_abandon_frame": 449,
      "bag_roi": [
        329,
        36,
        48,
        40
      ],
      "radius_px": 200,
      "threshold_s": 15
    },
	{
      "has_abandonment": true,
      "true_abandon_frame": 2012,
      "bag_roi": [
        244,
        23,
        55,
        11
      ],
      "radius_px": 200,
      "threshold_s": 15
    }
  ]
}
\end{lstlisting}

